\chapter{Lung Cancer}

\section*{Definition and Overview}
Lung cancer is a malignant epithelial neoplasm arising from the parenchyma or bronchial tree of the respiratory tract. It is histologically and clinically divided into two major categories: non-small cell lung cancer (NSCLC), which accounts for approximately 85\% of cases and encompasses adenocarcinoma, squamous cell carcinoma, and large cell carcinoma; and small cell lung cancer (SCLC), which constitutes roughly 15\% of cases and is characterized by a rapid cell doubling time, high growth fraction, and early haematogenous metastasis. Lung cancer is a highly aggressive disease and remains the leading cause of cancer-related mortality globally and in Australia, presenting a profound clinical and public health challenge.

\section*{Epidemiology}
Globally, lung cancer is the leading cause of oncology-related death, accounting for approximately 1.8 million deaths annually. In Australia, it is the fifth most commonly diagnosed cancer, with over 14,000 new cases diagnosed each year, but it remains the single largest cause of cancer-related death, accounting for approximately 18\% of all oncology deaths. The disease is primarily diagnosed in older adults, with the median age at diagnosis being 71 years. While historical incidence rates were significantly higher in males, gender patterns have shifted; male incidence rates are declining due to long-term reductions in tobacco consumption, while female incidence rates have plateaud or risen slightly. Crucially, a significant proportion of lung cancers occur in lifelong non-smokers, particularly among females and patients of East Asian descent.

\section*{Aetiology and Pathophysiology}
The development of lung cancer is driven by chronic exposure to environmental carcinogens and genetic susceptibility, leading to the accumulation of somatic mutations in pulmonary epithelial cells.
\begin{itemize}
    \item \textbf{Tobacco smoking:} The primary risk factor, accounting for approximately 85\% of all lung cancer cases. Tobacco smoke contains over 70 carcinogens, including polycyclic aromatic hydrocarbons and nitrosamines, which induce DNA adducts and promote oncogenic mutations.
    \item \textbf{Radon gas exposure:} An invisible, odourless, naturally occurring radioactive gas produced by the decay of uranium in soil and rocks. It is the second leading cause of lung cancer and the primary cause in non-smokers.
    \item \textbf{Asbestos exposure:} Workplace exposure to asbestos fibres significantly increases the risk of lung cancer. This risk acts synergistically with tobacco smoke, multiplying the overall risk of disease development.
    \item \textbf{Occupational carcinogens:} Exposure to industrial agents, including silica dust, arsenic, chromium, nickel, beryllium, and diesel exhaust fumes.
    \item \textbf{Ambient air pollution:} Chronic exposure to outdoor air pollution, specifically fine particulate matter (PM2.5), which is classified by the WHO as a Group 1 carcinogen.
    \item \textbf{Genetic driver alterations:} Oncogenic driver mutations in NSCLC, such as epidermal growth factor receptor (EGFR) mutations, anaplastic lymphoma kinase (ALK) or ROS1 rearrangements, and KRAS, BRAF, or RET mutations, which promote cell growth independent of external signals.
    \item \textbf{Tumour suppressor loss:} Near-universal inactivation of TP53 and RB1 genes, particularly in small cell lung cancer, disrupting normal cell cycle checkpoints.
    \item \textbf{Pre-existing pulmonary diseases:} Chronic obstructive pulmonary disease (COPD), idiopathic pulmonary fibrosis, and tuberculosis increase lung cancer risk independently of tobacco exposure.
\end{itemize}

\section*{Clinical Features}
The clinical features of lung cancer vary depending on tumour location, size, regional spread, and the presence of distant metastases or paraneoplastic phenomena.
\begin{itemize}
    \item \textbf{Persistent cough:} A new, worsening, or changing cough that does not resolve within three weeks is the most common presenting symptom.
    \item \textbf{Haemoptysis:} Coughing up blood or blood-streaked sputum, which requires immediate clinical investigation.
    \item \textbf{Dyspnea and wheezing:} Shortness of breath or localized wheezing caused by airway compression, post-obstructive atelectasis, or pleural effusion.
    \item \textbf{Chest pain:} Persistent, dull, or pleuritic chest pain indicating thoracic wall, rib, or pleural involvement.
    \item \textbf{Systemic B-symptoms:} Unexplained weight loss, cachexia, anorexia, and profound fatigue.
    \item \textbf{Voice hoarseness:} Caused by compression of the left recurrent laryngeal nerve by mediastinal lymph node metastases.
    \item \textbf{Superior vena cava (SVC) syndrome:} Facial swelling, neck vein distension, and upper extremity congestion caused by compression of the SVC by mediastinal masses.
    \item \textbf{Pancoast syndrome:} Apical lung tumours invading the brachial plexus and sympathetic chain, presenting as severe shoulder pain and Horner's syndrome (ptosis, miosis, anhidrosis).
    \item \textbf{Paraneoplastic syndromes:} Systemic manifestations, such as hypercalcaemia (from parathyroid hormone-related peptide secretion in squamous cell carcinoma) or syndrome of inappropriate antidiuretic hormone (SIADH) causing hyponatraemia in SCLC.
\end{itemize}

\section*{Differential Diagnosis}
Given its non-specific clinical and radiological presentation, lung cancer must be differentiated from several infectious and inflammatory conditions.
\begin{itemize}
    \item \textbf{Pulmonary tuberculosis:} Characterised by chronic cough, haemoptysis, night sweats, and apical cavitary lesions on chest imaging.
    \item \textbf{Pneumonia or lung abscess:} Acute or subacute respiratory infections presenting with cough, fever, purulent sputum, and pulmonary consolidation.
    \item \textbf{Systemic sarcoidosis:} A multi-system granulomatous disease that commonly presents with bilateral hilar lymphadenopathy and pulmonary nodules.
    \item \textbf{Pulmonary fungal infections:} Conditions like aspergillosis, histoplasmosis, or coccidioidomycosis, which can present as solitary or multiple pulmonary nodules.
    \item \textbf{Benign lung tumours:} Non-malignant masses, such as pulmonary hamartomas or bronchial adenomas.
    \item \textbf{Metastatic disease:} Malignant nodules in the lung originating from primary cancers of the colon, breast, kidney, or thyroid.
\end{itemize}

\section*{When to Seek Medical Care}
Patients experiencing a persistent cough lasting more than three weeks, chest pain, unexplained weight loss, or changes in a chronic respiratory cough should consult a general practitioner for further diagnostic workup, including a chest X-ray.

\textbf{Call 000 (or local emergency services) for:}
\begin{itemize}
    \item Haemoptysis (coughing up blood, especially in significant volumes).
    \item Acute, severe shortness of breath or inability to breathe.
    \item Sudden, severe, or crushing chest pain.
    \item Swelling of the face, neck, and upper body accompanied by difficulty breathing (indicative of superior vena cava obstruction).
    \item Sudden weakness, confusion, seizures, or loss of balance (potential signs of brain metastases).
\end{itemize}

\section*{Diagnosis and Investigations}
Establishing a diagnosis of lung cancer requires radiological staging and histopathological confirmation.
\begin{itemize}
    \item \textbf{Chest X-ray:} Frequently the first investigation, capable of identifying pulmonary masses, consolidated areas, or pleural effusions.
    \item \textbf{Computed Tomography (CT) of the chest and abdomen:} Highly detailed imaging to evaluate tumour size, location, relationship to local structures, and regional lymphadenopathy.
    \item \textbf{Positron Emission Tomography-CT (PET-CT):} Critical for clinical staging to identify regional lymph node involvement and distant metabolic metastases.
    \item \textbf{Tissue biopsy:} Performed via bronchoscopy (with endobronchial ultrasound-guided transbronchial needle aspiration, EBUS-TBNA) or CT-guided transthoracic core needle biopsy.
    \item \textbf{Molecular and biomarker profiling:} Routine testing of advanced NSCLC specimens for targetable alterations (EGFR, ALK, ROS1, BRAF, RET, NTRK) and PD-L1 expression to guide systemic therapy.
    \item \textbf{Brain Magnetic Resonance Imaging (MRI):} Indicated to rule out asymptomatic intracranial metastases in patients with diagnosed SCLC and stage II-IV NSCLC.
    \item \textbf{Pleural fluid aspiration (Thoracocentesis):} Cytological examination of fluid to confirm malignant pleural effusion in advanced stages.
\end{itemize}

\section*{Management}
Management strategies depend heavily on the histological type (NSCLC vs. SCLC), tumour stage, molecular driver status, and patient cardiorespiratory fitness.
\begin{itemize}
    \item \textbf{Surgical resection:} The gold standard curative treatment for early-stage (I-II) NSCLC, typically involving lobectomy or segmentectomy, often performed via video-assisted thoracic surgery (VATS).
    \item \textbf{Radiotherapy:} Stereotactic ablative radiotherapy (SABR) for medically inoperable early-stage NSCLC, or concurrent chemoradiotherapy for stage III disease.
    \item \textbf{Targeted molecular therapy:} Preferred first-line treatment for advanced NSCLC harbouring driver mutations, utilizing oral tyrosine kinase inhibitors such as osimertinib (EGFR) or alectinib (ALK).
    \item \textbf{Immunotherapy:} Immune checkpoint inhibitors (e.g.\ pembrolizumab, durvalumab) targeting PD-1, PD-L1, or CTLA-4, used alone or in combination with chemotherapy.
    \item \textbf{Chemotherapy:} Platinum-based doublet chemotherapy (cisplatin or carboplatin with gemcitabine, pemetrexed, or paclitaxel) remains the backbone for advanced NSCLC without targetable mutations and for SCLC.
    \item \textbf{Palliative and supportive care:} Early integration of symptom management for pain, dyspnea, and cough improves quality of life and prolongs survival in advanced disease.
\end{itemize}

\section*{Complications}
\begin{itemize}
    \item \textbf{Malignant pleural effusion:} Accumulation of fluid in the pleural cavity leading to lung collapse and severe dyspnea, often requiring indwelling pleural catheters or pleurodesis.
    \item \textbf{Post-obstructive pneumonia:} Bacterial infection arising in lung tissue distal to a bronchus obstructed by tumor growth.
    \item \textbf{Metastatic complications:} Pathological fractures and spinal cord compression from bone metastases; neurological deficits and seizures from brain metastases.
    \item \textbf{Venous thromboembolism (VTE):} Deep vein thrombosis and pulmonary embolism due to the highly prothrombotic nature of lung cancer.
    \item \textbf{Treatment-induced toxicities:} Radiation pneumonitis, immune-mediated colitis or pneumonitis, and severe chemotherapy-induced cytopenias.
\end{itemize}

\section*{Prognosis and Outcomes}
The overall prognosis for lung cancer remains challenging, but the advent of immunotherapies and targeted therapies has significantly improved survival outcomes. In Australia, the overall 5-year survival rate for lung cancer is approximately 20\% to 22\%. Stage at diagnosis remains the most critical prognostic factor; patients diagnosed with early, localised disease have a 5-year survival rate exceeding 60\%, whereas for those with distant metastatic disease, the 5-year survival is less than 5\%. Small cell lung cancer has a particularly aggressive course, with limited-stage disease having a median survival of 15 to 20 months, and extensive-stage disease having a median survival of 8 to 12 months.

\section*{Prevention and Self-Care}
\begin{itemize}
    \item \textbf{Smoking cessation:} The single most effective action to prevent lung cancer. The risk of developing lung cancer drops by half 10 years after quitting.
    \item \textbf{Avoidance of second-hand smoke:} Minimise exposure to passive environmental tobacco smoke in residential and workplace settings.
    \item \textbf{Radon mitigation:} Test indoor radon levels in homes located in high-risk regions and install mitigation systems if necessary.
    \item \textbf{Occupational safety compliance:} Strictly adhere to workplace safety guidelines, use respiratory protective equipment, and minimize exposure to asbestos, silica, and industrial chemicals.
    \item \textbf{Healthy dietary habits:} Consume a diet high in fresh fruits and vegetables, and maintain a regular physical exercise regime.
\end{itemize}

\section*{Sources and Further Reading}
The following reputable organisations provide evidence-based information, clinical guidelines, and support networks for lung cancer:
\begin{itemize}
    \item Lung Foundation Australia. \textit{Support, advocacy, and patient education resources for Australians affected by lung cancer.} https://lungfoundation.com.au/
    \item Cancer Council Australia. \textit{Evidence-based guidelines on lung cancer diagnosis, treatment, and support.} https://www.cancer.org.au/
    \item Thoracic Oncology Group of Australasia (TOGA). \textit{Clinical trial and research information on thoracic cancers.} https://thoraciconcology.org.au/
    \item National Institutes of Health (NIH) MedlinePlus. \textit{Patient-focused educational materials on lung cancer symptoms and causes.} https://medlineplus.gov/
    \item World Health Organization (WHO). \textit{Global cancer statistics, facts sheets, and tobacco control initiatives.} https://www.who.int/
\end{itemize}
